% Problem record op_61cc0e261b3889c6. This TeX file is derived from
% database/problems_json/op_61cc0e261b3889c6.json by scripts/sync-tex.mjs; edit the JSON record
% and rerun the script rather than editing this file.
\section{QMA(2) versus QMA}
\label{sec:op_61cc0e261b3889c6}

\paragraph{Problem.}
Is every promise problem verifiable by a quantum Merlin-Arthur protocol
with two unentangled witnesses also verifiable by a protocol with a single
arbitrary witness, that is, is $\mathsf{QMA}(2) = \mathsf{QMA}$ in the
standard, unrelativized setting?  A promise problem $L = (L_{yes}, L_{no})$
is in $\mathsf{QMA}(2)$ if some uniform polynomial-time quantum verifier,
given an input $x$ of length $n$ and two witnesses of $\operatorname{poly}(n)$
qubits each that are promised to be unentangled across the two registers,
accepts some product state $\sigma_1 \otimes \sigma_2$ with probability at
least $2/3$ for every $x \in L_{yes}$ and every product state
$\tau_1 \otimes \tau_2$ with probability at most $1/3$ for every
$x \in L_{no}$; $\mathsf{QMA}$ is the same model with a single,
unrestricted witness.  Discarding one witness gives
$\mathsf{QMA} \subseteq \mathsf{QMA}(2)$, and nondeterministically guessing
classical descriptions of the two witnesses and simulating the verifier
gives $\mathsf{QMA}(2) \subseteq \mathsf{NEXP}$, so
\begin{equation}
  \mathsf{QMA} \subseteq \mathsf{QMA}(2) \subseteq \mathsf{NEXP},
  \label{eq:qma2-containment}
\end{equation}
while the SWAP-test reduction of Harrow and Montanaro gives
$\mathsf{QMA}(k) = \mathsf{QMA}(2)$ for every constant $k \ge 2$.  The open
question is whether the first inclusion in Eq.~\eqref{eq:qma2-containment}
is an equality: can every protocol whose soundness is required only against
separable pairs of witnesses be simulated by a single-witness protocol with
polynomial overhead?  No improvement of either containment in
Eq.~\eqref{eq:qma2-containment} is known in the unrelativized setting, and
relativized evidence such as an oracle separation does not settle this
question.

\subsection*{Status}

Unsolved

\subsection*{Source}

The classes of problems verifiable with multiple unentangled quantum
certificates were introduced as $\mathsf{QMA}(k)$ by Kobayashi, Matsumoto,
and Yamakami \sourcecite{ref:qma2-kmy01}{KMY01}, and whether the
unentanglement promise adds verification power has been open since; Jeronimo
and Wu state that for $\mathsf{QMA}(2)$ itself only the trivial containments
of Eq.~\eqref{eq:qma2-containment} are known
\sourcecite{ref:qma2-jw24}{JW24}, as do Bostanci et al.
\sourcecite{ref:qma2-bgh26}{BGH+26}.

\subsection*{Progress}

\begin{itemize}
  \item Kobayashi, Matsumoto, and Yamakami defined $\mathsf{QMA}(k)$, quantum
  certificate verification with $k$ unentangled certificates, in a work
  devoted to the single-versus-multiple comparison
  \sourcecite{ref:qma2-kmy01}{KMY01}.

  \item Harrow and Montanaro proved $\mathsf{QMA}(k) = \mathsf{QMA}(2)$ for every
  $k \ge 2$ by a SWAP-test reduction, so two unentangled witnesses already
  capture any constant number of them \sourcecite{ref:qma2-hm13}{HM13}.

  \item For the non-negative-amplitude variant, Jeronimo and Wu proved
  $\mathsf{QMA}^{+}(2) = \mathsf{NEXP}$ and showed that $\mathsf{QMA}(2)$ with
  a sufficiently large constant completeness-soundness gap coincides with
  $\mathsf{QMA}^{+}(2)$, while stating that for $\mathsf{QMA}(2)$ itself only
  the trivial containments of Eq.~\eqref{eq:qma2-containment} are known
  \sourcecite{ref:qma2-jw24}{JW24}.

  \item Bassirian, Fefferman, Leigh, Marwaha, and Wu separated a single witness
  from two unentangled witnesses only for the modified class
  $\mathsf{QMA}_{IS}$ of internally separable proofs and only assuming
  $\mathsf{EXP} \ne \mathsf{NEXP}$, framing $\mathsf{QMA}(2) = \mathsf{NEXP}$ as
  the open target and leaving $\mathsf{QMA}(2)$ versus $\mathsf{QMA}$
  untouched \sourcecite{ref:qma2-bfl24}{BFL+24}.

  \item The survey of Jeronimo, Wu, and Leigh is devoted to the
  $\mathsf{QMA}(2)$ universe of complexity, entanglement, and optimization
  questions and describes the class as central yet poorly understood
  \sourcecite{ref:qma2-jwl26}{JWL26}.

  \item Bostanci, Grewal, Haferkamp, Huang, Hwang, Natarajan, and Nirkhe
  constructed the first quantum, that is unitary, oracle relative to which
  $\mathsf{QMA} \ne \mathsf{QMA}(2)$, resolving Watrous's no-disentanglers
  conjecture that every $(\varepsilon, \delta)$-disentangler with
  $\varepsilon + \delta < 1$ requires input size exponential in the number of output qubits.  They
  state that in the unrelativized setting essentially all that is known is
  the chain of Eq.~\eqref{eq:qma2-containment}, that a classical oracle
  separation still eludes us, and that any in-place amplification implying
  $\mathsf{QMA}(2) \subseteq \mathsf{QMA}$ must be non-relativizing; the
  preprint, submitted 2 September 2026, is not yet refereed
  \sourcecite{ref:qma2-bgh26}{BGH+26}.
\end{itemize}

\subsection*{References}

\begin{enumerate}[label={},leftmargin=4.8em,labelsep=0.5em]
  \footnotesize
  \item[\textup{[KMY01]}]\label{ref:qma2-kmy01}
  H. Kobayashi, K. Matsumoto, and T. Yamakami, ``Quantum certificate
  verification: Single versus multiple quantum certificates,''
  \href{https://arxiv.org/abs/quant-ph/0110006}{arXiv:quant-ph/0110006}
  (2001).

  \item[\textup{[HM13]}]\label{ref:qma2-hm13}
  A. W. Harrow and A. Montanaro, ``Testing Product States, Quantum
  Merlin-Arthur Games and Tensor Optimization,'' \emph{Journal of the ACM}
  \textbf{60}(1), Article 3 (2013).
  \href{https://doi.org/10.1145/2432622.2432625}{doi:10.1145/2432622.2432625}.

  \item[\textup{[JW24]}]\label{ref:qma2-jw24}
  F. G. Jeronimo and P. Wu, ``The Power of Unentangled Quantum Proofs with
  Non-negative Amplitudes,''
  \href{https://arxiv.org/abs/2402.18790}{arXiv:2402.18790} (2024).

  \item[\textup{[BFL+24]}]\label{ref:qma2-bfl24}
  R. Bassirian, B. Fefferman, I. Leigh, K. Marwaha, and P. Wu, ``Quantum
  Merlin-Arthur with an internally separable proof,''
  \href{https://arxiv.org/abs/2410.19152}{arXiv:2410.19152} (2024).

  \item[\textup{[JWL26]}]\label{ref:qma2-jwl26}
  F. G. Jeronimo, P. Wu, and I. Leigh, ``The QMA(2) Universe: Complexity,
  Entanglement, and Optimization,'' \emph{ACM SIGACT News} \textbf{57}(1),
  64--99 (2026).
  \href{https://doi.org/10.1145/3802807.3802814}{doi:10.1145/3802807.3802814}.

  \item[\textup{[BGH+26]}]\label{ref:qma2-bgh26}
  J. Bostanci, S. Grewal, J. Haferkamp, A. Huang, Y. Hwang, A. Natarajan,
  and C. Nirkhe, ``A quantum oracle separation between QMA(2) and QMA,''
  \href{https://arxiv.org/abs/2609.02865}{arXiv:2609.02865} (2026).
\end{enumerate}

\subsection*{Comment}

What remains unresolved is the unrelativized class equality: no improvement
of either containment in Eq.~\eqref{eq:qma2-containment} is known.  A proof
that $\mathsf{QMA}(2) \subseteq \mathsf{QMA}$ settles it with equality; the
quantum-oracle separation of \sourcecite{ref:qma2-bgh26}{BGH+26} shows that
such a simulation must use non-relativizing techniques, and its resolution
of the no-disentanglers conjecture closes the query-efficient route.
Conversely, exhibiting any language in
$\mathsf{QMA}(2) \setminus \mathsf{QMA}$ settles it with a strict
inclusion; in particular, proving $\mathsf{QMA}(2) = \mathsf{NEXP}$, the
target of the non-negative-amplitude and internally-separable-proof
programs, would separate the classes assuming
$\mathsf{EXP} \ne \mathsf{NEXP}$, since $\mathsf{QMA} \subseteq \mathsf{PP}
\subseteq \mathsf{PSPACE} \subseteq \mathsf{EXP}$.  Weaker decisive
progress would improve a single containment, for example
$\mathsf{QMA}(2) \subseteq \mathsf{EXP}$ or
$\mathsf{NEXP} \subseteq \mathsf{QMA}(2)$ at any constant
completeness-soundness gap.  Finding a classical oracle separating
$\mathsf{QMA}$ from $\mathsf{QMA}(2)$ is a further open relativized
sub-question, distinct from the unitary-oracle version resolved in 2026.
The zoo's neighboring records pose different questions: ``The quantum PCP
conjecture'' asks for QMA-hardness of constant-relative-gap local
Hamiltonians on a promise problem, and ``Constant trace-distance
separability testing'' relates its promise gap to estimating
$\mathsf{QMA}(2)$ acceptance probabilities.

\subsection*{Field}

Quantum algorithm

\subsection*{Topic}

Computational complexity and computability; Quantum separability

\subsection*{ID}

\texttt{op\_61cc0e261b3889c6}
